\documentclass[11pt,letterpaper]{article}
\usepackage[margin=1in]{geometry}
\usepackage{amsmath,amssymb,amsthm,booktabs,tabularx,hyperref,microtype}
\hypersetup{colorlinks=true,linkcolor=blue,citecolor=blue,urlcolor=blue}
\newtheorem{proposition}{Proposition}
\newtheorem{corollary}[proposition]{Corollary}
\newcommand{\CE}{\operatorname{CE}}
\newcommand{\KL}{\operatorname{KL}}
\newcommand{\TV}{\operatorname{TVD}}
\newcommand{\Var}{\operatorname{Var}}
\newcommand{\Cov}{\operatorname{Cov}}
\newcommand{\E}{\mathbb{E}}
\newcommand{\R}{\mathbb{R}}
\DeclareMathOperator{\softmax}{softmax}
\title{What Correlation Can and Cannot Identify in Human-Rating Judges}
\author{Brando Miranda\\Stanford University\\\texttt{brando9@stanford.edu}}
\date{11-23-2026}
\begin{document}
\maketitle
\begin{abstract}
% CS197 move 1: problem motivation (failure mode) -- perfect association can hide a wrong rating distribution.
A judge can achieve perfect Pearson correlation with human mean ratings while predicting the wrong conditional rating distribution.
% CS197 move 2: assumption (setup bit) -- association does not itself imply distribution recovery.
Interpreting this agreement requires separating identification of a mean from identification of a distribution.
% CS197 move 3: contribution (bit flip) -- characterize the objective rather than claim empirical training superiority.
We characterize exactly what correlation of conditional means identifies on a finite ordered rating scale.
% CS197 move 4a: instantiation (bit flip instantiation) -- precise invariances and feasible model directions.
The objective depends only on means and is invariant to a common positive affine transformation.
% CS197 move 4b: instantiation (bit flip instantiation) -- state the model-class qualifications.
We derive its gradients, characterize restricted-family identification, and distinguish binary ratings from scales with independent shape variation.
% CS197 move 4c: instantiation (bit flip instantiation) -- compare population targets using self-contained proofs.
Cross-entropy and normalized individual-rating ranked probability score identify the distribution in expectation; their finite-data training performance requires separate analysis.
% CS197 move 4d: instantiation (bit flip instantiation) -- name the exact loss decompositions.
We also derive a signed cross-entropy decomposition and an exact cumulative-probability decomposition of excess ranked probability score.
% CS197 move 4e: instantiation (bit flip instantiation) -- distinguish information bounds from an estimator guarantee.
Recovery lower bounds show why exact means alone cannot determine distributional shape without additional restrictions.
% CS197 move 5: evaluation (diagnostic result) -- exact rational arithmetic in experiments/exact_counterexample/reproduce.py, no trained models.
An exact three-item example has correlation one, average total variation distance $1/5$, and normalized excess ranked probability score $3/200$.
% CS197 move 6: implication (north star) -- align evaluation with its target and limit finite-data claims.
The analysis separates association, distribution recovery, and estimation error, providing conditions for meaningful comparisons without asserting a universal advantage for cross-entropy training.
\end{abstract}

\section{Introduction}
Perfect correlation with human mean ratings does not identify a judge's rating distribution.
A conditional distribution describes how human ratings vary for an item.
Its mean summarizes only one aspect of that variation.
If two forecasts have identical means on every item, correlation of those means assigns identical values to both forecasts.
The distributions can still disagree about uncertainty and the probabilities of individual ratings.

The interpretation of correlation therefore depends on the prediction target.
When the target is association between scalar scores, correlation directly describes that association.
When the target is a conditional rating distribution, association alone leaves additional questions unanswered.
Strictly proper scoring rules address distribution recovery by making the true distribution the unique optimum of expected score~\cite{gneiting2007}.
The ranked probability score uses the ordering of the rating categories~\cite{epstein1969}.
Our comparison concerns these mathematical targets rather than a claim about how often practitioners use any objective.

We characterize the information discarded by Pearson correlation of conditional means.
At correlation one, forecast means equal target means after a common positive affine transformation.
On a scale with at least three categories, unrestricted distributions also retain shape variation at fixed means.
This result concerns the objective's dependence on outputs.
It places no linearity restriction on a neural judge's map from an item to its mean rating.

The model class determines which ambiguities are feasible.
A restricted family can tie shape to the mean or exclude every alternative affine copy of the target means.
Shared parameters can also couple changes in mean and shape.
We account for these restrictions through the model's mean map and its feasible directions.
We then compare population cross-entropy, normalized individual-rating ranked probability score, and total variation distance to the true distribution.

We give detailed proofs and an exact finite illustration rather than a trained-model comparison.
The illustration holds the means fixed across three strictly positive rating distributions.
Its correlation is one, although its average normalized excess ranked probability score is $3/200$.
The supplied arithmetic reproduces these values without fitting a model.
A separate analysis explains why finite rating noise and estimation error prevent a universal training comparison.

The results support matching the evaluation criterion to the intended output.
They do not establish that one training algorithm always generalizes better.
We make the following contributions:
\begin{enumerate}
\item We characterize the invariances, equality conditions, gradients, and probability-profile variants of Pearson correlation.
\item We derive distributional loss identities, moment bounds, and lower bounds for recovery from means alone.
\item We state how the rating scale, model restrictions, and finite-data assumptions change these conclusions.
\end{enumerate}

\section{Problem setup and notation}
Fix $n\ge2$ items $x_1,\ldots,x_n$ and ordered ratings $R=\{0,\ldots,K-1\}$ with $K\ge2$.
For item $i$, let $P_i$ denote the true conditional human-rating distribution and let $q_i$ denote a forecast.
Both belong to the probability simplex $\Delta^{K-1}$.
Define
\[
 t_i=\sum_{r=0}^{K-1}rP_i(r),\qquad
 m_i=\sum_{r=0}^{K-1}rq_i(r),\qquad
 \bar t=\frac1n\sum_i t_i,\qquad \bar m=\frac1n\sum_i m_i.
\]
An item can include a prompt, reference answer, and candidate answer; all conditioning information is included in $x_i$.
A model class is $\mathcal Q=\{q_\theta:\theta\in\Theta\}$ with mean map $M(\theta)=(m_\theta(x_1),\ldots,m_\theta(x_n))$.
Identification here concerns these conditional distributions on the specified items unless a population of inputs is explicitly introduced.
Different parameters may represent the same predictions.

Write $s_t^2=n^{-1}\sum_i(t_i-\bar t)^2$ and $s_m^2=n^{-1}\sum_i(m_i-\bar m)^2$.
For $s_t,s_m>0$, Pearson correlation across item means is
\begin{equation}\label{eq:pearson}
 \rho(m,t)=\frac{\sum_i(m_i-\bar m)(t_i-\bar t)}{n s_m s_t}.
\end{equation}
If either vector is constant, this coefficient is undefined.
The minimized correlation loss is $1-\rho$.
Whenever this paper says ``correlation'' without a qualifier, it means Eq.~\eqref{eq:pearson}.

For distributions $P,q$, use natural logarithms and define entropy, cross-entropy (CE), Kullback--Leibler divergence (KL), and total variation distance (TVD) by
\begin{align}
 H(P)&=-\sum_r P(r)\log P(r), &
 \CE(P,q)&=-\sum_r P(r)\log q(r),\\
 \KL(P\Vert q)&=\sum_r P(r)\log\frac{P(r)}{q(r)}, &
 \TV(P,q)&=\frac12\sum_r|P(r)-q(r)|.
\end{align}
Terms with $P(r)=0$ contribute zero.
If $P(r)>0=q(r)$, CE and KL equal $+\infty$.
Finite CE requires $q$ positive on the support of $P$.
Some decompositions below require $q(r)>0$ for all $r$ and state that stronger condition explicitly.
Entropy, CE, and KL use nats; TVD lies in $[0,1]$.

Let $F_q(k)=\sum_{r=0}^k q(r)$ for $k=0,\ldots,K-2$.
The normalized individual-rating ranked probability score (RPS) is
\begin{equation}\label{eq:rps}
 S(q,y)=\frac1{K-1}\sum_{k=0}^{K-2}\big(F_q(k)-\mathbf1\{y\le k\}\big)^2.
\end{equation}
This score lies in $[0,1]$, and smaller values are better.
We use $\mathcal E_P(q)=\E_{Y\sim P}S(q,Y)-\E_{Y\sim P}S(P,Y)$ for its excess expected score.
The original ranked-category construction and the general proper-scoring framework appear in Epstein~\cite{epstein1969} and Gneiting and Raftery~\cite{gneiting2007}, respectively.
All identities needed here are proved below.
Unless stated otherwise, losses and distances across items are averaged with weight $1/n$.
An empirical histogram estimates $P_i$ from finitely many ratings; it is not assumed to equal $P_i$.

\section{What correlation identifies}
\begin{proposition}[Dependence on means and affine invariance]\label{prop:invariance}
For nonconstant $m,t$,
\begin{equation}\label{eq:meanonly}
 \rho(m,t)=\frac{\sum_i(t_i-\bar t)\sum_r r q_i(r)}{\sqrt{\sum_i(t_i-\bar t)^2\sum_i(m_i-\bar m)^2}}.
\end{equation}
Thus forecasts with equal item means have equal correlation.
For any $a\in\R$ and $b>0$, $\rho(a\mathbf1+bm,t)=\rho(m,t)$.
Moreover,
\begin{equation}\label{eq:perfect}
 \rho(m,t)=1\quad\Longleftrightarrow\quad m=a\mathbf1+bt\quad\text{for some }a\in\R,\ b>0.
\end{equation}
The right side must be feasible in the model class to describe an attainable optimum.
\end{proposition}
\begin{proof}
Since $\sum_i(t_i-\bar t)=0$,
\[
 \sum_i(m_i-\bar m)(t_i-\bar t)=\sum_i m_i(t_i-\bar t).
\]
Substitution of $m_i=\sum_r r q_i(r)$ gives Eq.~\eqref{eq:meanonly}.
Its numerator is linear in the probabilities; its denominator depends on the forecast only through $m$.
The centered vector of $a\mathbf1+bm$ equals $b(m-\bar m\mathbf1)$.
Its covariance numerator and standard-deviation denominator both acquire the same factor $b$, proving invariance.
Cauchy--Schwarz bounds the correlation by one.
Equality with positive sign holds exactly when $m-\bar m\mathbf1=b(t-\bar t\mathbf1)$ for $b>0$.
Setting $a=\bar m-b\bar t$ proves both directions of Eq.~\eqref{eq:perfect}.
\end{proof}

The linear functional from probabilities to a mean does not constrain the map $x\mapsto m_\theta(x)$ to be linear.
A nonlinear mean predictor can attain correlation one whenever it can represent the nonconstant target mean function.
The objective still cannot distinguish feasible alternatives with identical means.

\begin{proposition}[Rating scale and dimension]\label{prop:dimension}
For binary ratings, $q=(1-m,m)$ is determined by its mean.
For $K\ge3$, the set of distributions with a fixed interior mean $m\in(0,K-1)$ has affine dimension $K-2$.
With unrestricted item distributions, the all-positive part of the correlation-one set has dimension $n(K-2)+2$.
Boundary constraints or model restrictions can reduce the dimension.
\end{proposition}
\begin{proof}
For $K=2$, $m=q(1)$ and mass conservation gives $q(0)=1-m$.
For general $K$, the constraints $\sum_r q(r)=1$ and $\sum_r r q(r)=m$ have independent coefficient vectors.
Their common affine solution space therefore has dimension $K-2$.
Every interior mean admits an all-positive distribution, for example the exponential family established in Proposition~\ref{prop:exponential}.
At such a point, positivity imposes no additional local equality constraints.
For $n$ items with independently specified means, the dimensions add to $n(K-2)$.
In Eq.~\eqref{eq:perfect}, the columns $\mathbf1$ and $t$ are independent because $t$ is nonconstant.
The feasible all-positive stratum allows two locally independent changes in $(a,b)$, giving the stated total dimension.
\end{proof}

For $K=3$ and mean one, write $q_s=(s,1-2s,s)$ with $0\le s\le1/2$.
Its variance is $2s$, although every distribution in the family has the same mean.
For binary targets with means $t_i$ and feasible correlation-one forecasts $m_i=a+bt_i$, the excess expected RPS is
\[
 \frac1n\sum_i\big(a+(b-1)t_i\big)^2.
\]
Its excess CE is $n^{-1}\sum_i\KL(\operatorname{Bern}(t_i)\Vert\operatorname{Bern}(a+bt_i))$.
Thus binary ratings remove independent shape ambiguity but retain level and scale ambiguity.
Here RPS uses the one-threshold binary Brier convention; the sum of squared errors over both category probabilities would be twice as large.

\begin{proposition}[Gradients and feasible directions]\label{prop:gradient}
Let $q_i=\softmax(u_i)$ and let $z_i=(m_i-\bar m)/s_m$ and $\zeta_i=(t_i-\bar t)/s_t$.
With $g_i=(\zeta_i-\rho z_i)/(n s_m)$,
\begin{equation}\label{eq:gradients}
 \frac{\partial\rho}{\partial m_i}=g_i,\qquad
 \frac{\partial\rho}{\partial u_i(r)}=g_i q_i(r)(r-m_i),\qquad
 \frac{\partial\sum_j\CE(P_j,q_j)}{\partial u_i(r)}=q_i(r)-P_i(r).
\end{equation}
For a differentiable shared-parameter model, $D\rho(\theta)[v]=0$ whenever $DM(\theta)[v]=0$.
Along a feasible path whose means are exactly constant, correlation is exactly constant wherever defined.
\end{proposition}
\begin{proof}
Let $T=t-\bar t\mathbf1$, $C=m-\bar m\mathbf1$, and $N=T^\top m$.
Then $\rho=N/(\|T\|\|C\|)$, $\partial N/\partial m_i=T_i$, and $\partial\|C\|/\partial m_i=C_i/\|C\|$ because $\sum_i C_i=0$.
The quotient rule gives
\[
 \frac{\partial\rho}{\partial m_i}=\frac{T_i}{\|T\|\|C\|}-\rho\frac{C_i}{\|C\|^2}=\frac{\zeta_i-\rho z_i}{n s_m}.
\]
The softmax derivative is $\partial q_i(s)/\partial u_i(r)=q_i(s)(\mathbf1\{s=r\}-q_i(r))$.
Consequently,
\[
 \frac{\partial m_i}{\partial u_i(r)}=\sum_s s q_i(s)(\mathbf1\{s=r\}-q_i(r))=q_i(r)(r-m_i).
\]
The chain rule gives the correlation derivative.
Differentiating $-\sum_sP_i(s)\log q_i(s)$ gives $-P_i(r)+q_i(r)\sum_sP_i(s)=q_i(r)-P_i(r)$.
The shared-parameter derivative is $\sum_i g_i DM_i(\theta)[v]$, which vanishes under the stated condition.
Exact constancy along a mean-preserving path follows directly from Proposition~\ref{prop:invariance}.
\end{proof}

For the minimized loss $1-\rho$, the correlation derivative in Eq.~\eqref{eq:gradients} changes sign.
Averaging the summed CE divides its derivative by $n$.
The ambient probability derivative can be written $\partial\rho/\partial q_i(r)=r g_i$; feasible simplex derivatives are its restrictions to mass-preserving directions.
An independently parameterized shape output that never changes $M$ receives zero direct gradient from correlation.
A shared parameter can change both mean and shape, so correlation can change shape indirectly.
A first-order mean-preserving direction guarantees a zero directional derivative at that point, not constancy at higher orders.

\begin{proposition}[Correlation of probability profiles]\label{prop:profiles}
Let $U=(1/K,\ldots,1/K)$ and let $P$ be nonuniform.
For any $b>0$ such that $q=U+b(P-U)$ is feasible, Pearson correlation across the $K$ category probabilities equals one.
In particular, all $0<b\le1$ are feasible.
The same construction with a common $b$ across all items yields correlation one after stacking every item/category pair, provided the stacked target is nonconstant.
\end{proposition}
\begin{proof}
The average entry of each probability profile is $1/K$.
Its centered vectors therefore satisfy $q-U=b(P-U)$, and Proposition~\ref{prop:invariance}'s equality argument applies.
For $0<b\le1$, $q=(1-b)U+bP$ is a probability distribution.
The stacked vectors also have average entry $1/K$ and satisfy the same proportionality when $b$ is common.
\end{proof}

Itemwise profile correlations can use different $b_i$; their average is a different objective from correlation of stacked profiles.
Uniform targets or forecasts make the corresponding profile correlation undefined.
These objectives can distinguish some equal-mean shapes because they use probability contrasts.
They still fail to identify the distribution uniquely, and they are distinct from Eq.~\eqref{eq:pearson}.

\section{Distributional losses and recovery bounds}
\subsection{Cross-entropy and the signed shape decomposition}
\begin{proposition}[Strict propriety of CE]\label{prop:ce}
For every true distribution $P$,
\[
 \CE(P,q)=H(P)+\KL(P\Vert q),\qquad \KL(P\Vert q)\ge0,
\]
with equality in the second relation exactly when $q=P$.
Thus expected individual-rating CE uniquely identifies $P$ over the unrestricted simplex.
\end{proposition}
\begin{proof}
Subtracting $H(P)$ from CE gives the displayed KL expression whenever CE is finite; the identity extends to infinite CE.
For $A=\{r:P(r)>0\}$ and a forecast positive on $A$, concavity of the logarithm gives
\[
 -\KL(P\Vert q)=\sum_{r\in A}P(r)\log\frac{q(r)}{P(r)}
 \le\log\sum_{r\in A}q(r)\le0.
\]
Equality requires $q(r)/P(r)$ constant on $A$ and $\sum_{r\in A}q(r)=1$, forcing that constant to be one.
If $q$ is zero on a positive-probability category, KL is infinite and cannot minimize CE.
\end{proof}

\begin{proposition}[Natural exponential families and moment matching]\label{prop:exponential}
Let $h(r)>0$ on $R$ and
\[
 q_\eta(r)=h(r)e^{\eta r-A(\eta)},\qquad
 A(\eta)=\log\sum_r h(r)e^{\eta r},\qquad \eta\in\R.
\]
Then $A'(\eta)=m(\eta)$ and $A''(\eta)=\Var_{q_\eta}(R)>0$.
The mean map is a bijection from $\R$ to $(0,K-1)$.
For any target $P$ with interior mean $t$, expected CE within this family has a unique finite minimizer satisfying $m(\eta)=t$, even if $P$ lies outside the family.
For $h\equiv1$, this mean-matching distribution, denoted $q_t$, uniquely maximizes entropy among all distributions with mean $t$.
\end{proposition}
\begin{proof}
Differentiating the finite sum defining $A$ gives $A'=\sum_r r q_\eta(r)$.
Differentiating again gives $A''=\sum_r r^2q_\eta(r)-(\sum_r rq_\eta(r))^2$.
All categories have positive probability and at least two distinct ratings exist, so this variance is strictly positive.
Hence $m$ is continuous and strictly increasing.
As $\eta\to-\infty$, mass concentrates at zero, and as $\eta\to+\infty$, mass concentrates at $K-1$.
The intermediate value theorem establishes the claimed range and inverse.
Writing $t=\sum_r rP(r)$ gives
\[
 \CE(P,q_\eta)=-\sum_rP(r)\log h(r)-\eta t+A(\eta).
\]
Its first derivative is $m(\eta)-t$ and its second derivative is strictly positive.
The unique zero of the first derivative is therefore the unique global minimizer.
For $h\equiv1$ and $\eta$ chosen to give mean $t$, every distribution $P$ with that mean satisfies
\[
 \KL(P\Vert q_t)=-H(P)-\eta t+A(\eta)=H(q_t)-H(P).
\]
Nonnegativity and the equality condition of KL prove unique entropy maximization.
\end{proof}

At an endpoint target mean, $P$ is the corresponding point mass.
A full-support exponential family approaches that optimum as $\eta\to\pm\infty$; it does not attain it at a finite parameter.
Moment matching here requires the sufficient statistic $r$ and the positive-variance condition.
It is not a property of every one-parameter rating family.
With freely chosen $\eta_i$ per item and interior feasible means, correlation-one predictions remain $m_i=a+bt_i$.
The family's shape is tied to these means, so the remaining ambiguity concerns level and scale.

Let $\eta_t$ and $\eta_m$ have nearby interior means $t$ and $m$, and write $V(t)=A''(\eta_t)>0$.
Then
\begin{align}
 \KL(q_{\eta_t}\Vert q_{\eta_m})
 &=A(\eta_m)-A(\eta_t)-A'(\eta_t)(\eta_m-\eta_t)\\
 &=\frac{(m-t)^2}{2V(t)}+o((m-t)^2).\label{eq:localprice}
\end{align}
Indeed, Taylor expansion gives $\tfrac12 V(t)(\eta_m-\eta_t)^2+o((\eta_m-\eta_t)^2)$, and the inverse-function derivative gives $\eta_m-\eta_t=(m-t)/V(t)+o(m-t)$.
For a misspecified target of mean $t$, subtracting its CE at $\eta_t$ from its CE at $\eta_m$ gives the same exact expression on the first line.
This is a local cost of mean miscalibration inside this family.
It does not assert that the family's variance equals the target variance.

\begin{proposition}[Signed CE decomposition]\label{prop:cross}
Let $P$ have interior mean $t$ and let $q_t$ be its maximum-entropy mean-matching distribution from Proposition~\ref{prop:exponential} with $h\equiv1$.
For a forecast $q$ positive on every category, define
\[
 G(P)=H(q_t)-H(P)\ge0,\qquad
 C(P,q)=\sum_r(P(r)-q_t(r))(-\log q(r)).
\]
Then
\begin{equation}\label{eq:cross}
 \CE(P,q)=H(P)+G(P)+\KL(q_t\Vert q)+C(P,q).
\end{equation}
The cross term $C(P,q)$ can be negative.
If $\log q(r)=\alpha+\beta r$, it vanishes, and the minimum excess CE within that log-affine-in-rating family is $G(P)$.
\end{proposition}
\begin{proof}
Add and subtract $\CE(q_t,q)$ to obtain $\CE(P,q)=\CE(q_t,q)+C(P,q)$.
Proposition~\ref{prop:ce} gives $\CE(q_t,q)=H(q_t)+\KL(q_t\Vert q)$.
Substitution of $H(q_t)=H(P)+G(P)$ yields Eq.~\eqref{eq:cross}.
If $\log q(r)=\alpha+\beta r$, then
\[
 C(P,q)=-\alpha\sum_r(P(r)-q_t(r))-\beta\sum_r r(P(r)-q_t(r))=0
\]
because both total mass and mean agree.
KL is uniquely minimized at $q=q_t$, establishing the family-specific excess $G(P)$.
For any positive $P$, substituting the unrestricted forecast $q=P$ into Eq.~\eqref{eq:cross} yields
\[
 C(P,P)=-G(P)-\KL(q_t\Vert P).
\]
This is negative whenever $P\ne q_t$.
\end{proof}

The nonnegative shape gap is therefore a restriction cost for the stated log-affine family, not an unavoidable penalty for every forecast.
The cross term prevents interpreting Eq.~\eqref{eq:cross} as a sum of nonnegative errors.
Other scalar-to-distribution families require their own analysis.
At endpoint means, the target and its maximum-entropy representative coincide and $G(P)=C(P,q)=0$.

\subsection{Ranked probability score and distribution distances}
\begin{proposition}[Expected RPS, signed CDF identity, and decomposition]\label{prop:rps}
For any $P,q$, let $d_k=F_P(k)-F_q(k)$, $\bar d=(K-1)^{-1}\sum_k d_k$, and $D(P,q)=\sum_k d_k^2$.
Then
\begin{align}
 \E_P S(q,Y)&=\frac1{K-1}\sum_k F_P(k)(1-F_P(k))+\frac{D(P,q)}{K-1},\label{eq:riskrps}\\
 \sum_{k=0}^{K-2}d_k&=m_q-m_P,\label{eq:signedcdf}\\
 \mathcal E_P(q)&=\frac{(m_q-m_P)^2}{(K-1)^2}+\frac1{K-1}\sum_k(d_k-\bar d)^2.\label{eq:rpssplit}
\end{align}
Consequently, individual-rating normalized RPS is strictly proper.
Also $|m_q-m_P|\le(K-1)\TV(P,q)$.
\end{proposition}
\begin{proof}
For $Z_k=\mathbf1\{Y\le k\}$, $\E Z_k=F_P(k)$ and $\Var(Z_k)=F_P(k)(1-F_P(k))$.
The identity $\E(a-Z_k)^2=(a-\E Z_k)^2+\Var(Z_k)$, summed and divided by $K-1$, proves Eq.~\eqref{eq:riskrps}.
Interchanging the finite sums gives
\[
 \sum_{k=0}^{K-2}F_P(k)=\sum_rP(r)(K-1-r)=K-1-m_P.
\]
Subtracting its counterpart for $q$ proves Eq.~\eqref{eq:signedcdf}, including its sign.
The centered sum-of-squares identity is
\[
 \sum_kd_k^2=\frac{(\sum_kd_k)^2}{K-1}+\sum_k(d_k-\bar d)^2.
\]
Dividing by $K-1$ gives Eq.~\eqref{eq:rpssplit}.
The excess vanishes exactly when all CDF values agree, which determines every category probability by successive differences.
This proves strict propriety, including distributions on the simplex boundary.
Finally, mass conservation gives
\[
 |m_q-m_P|=\left|\sum_r\left(r-\frac{K-1}{2}\right)(q(r)-P(r))\right|
 \le\frac{K-1}{2}\sum_r|q(r)-P(r)|.
\]
\end{proof}

The first term of Eq.~\eqref{eq:riskrps} is the irreducible expected score for a fresh stochastic rating.
It is distinct from excess RPS, which can be zero.
Equation~\eqref{eq:rpssplit} has mean-error denominator $(K-1)^2$, which is $25$ on a six-category scale.
The remaining term measures variation of the CDF differences after their average is removed.
It vanishes exactly when all CDF differences are equal; it is not an independent variance error.
If means agree, all remaining CDF discrepancy appears in this term.

TVD to the true distribution equals zero exactly when $q=P$.
It is a distribution distance rather than a strictly proper individual-rating score.
In fact,
\[
 \TV(q,\delta_y)=\frac12\left(1-q(y)+\sum_{r\ne y}q(r)\right)=1-q(y),
\]
so $\E_{Y\sim P}\TV(q,\delta_Y)=1-\sum_rP(r)q(r)$.
This expected loss is minimized by forecasts supported on the modal categories of $P$.
It generally does not recover $P$.
TVD to a sparse empirical histogram is also a noisy diagnostic, whose expectation need not equal TVD to $P$.

\begin{proposition}[CE controls bounded moments]\label{prop:pinsker}
For every $f:R\to[0,1]$,
\begin{equation}\label{eq:momentbound}
 |\E_P f-\E_q f|\le\TV(P,q)\le\sqrt{\KL(P\Vert q)/2}.
\end{equation}
In particular, the bound applies to every normalized raw moment $f(r)=(r/(K-1))^j$, $j\ge1$.
\end{proposition}
\begin{proof}
Write $\delta(r)=P(r)-q(r)$ and $A=\{r:\delta(r)\ge0\}$.
Since $\sum_r\delta(r)=0$, $\sum_{r\in A}\delta(r)=\TV(P,q)=d$.
The positive and negative parts bound $\sum_r f(r)\delta(r)$ above by $d$ and below by $-d$.
To prove the second inequality, suppose KL is finite; otherwise it is immediate.
Grouping categories in $A$ and its complement cannot increase KL.
For completeness, for any group $B$, concavity of the logarithm applied with weights $P(r)/P(B)$ yields
\[
 \sum_{r\in B}P(r)\log\frac{P(r)}{q(r)}\ge P(B)\log\frac{P(B)}{q(B)}.
\]
Zero group masses follow by continuity.
Set $a=P(A)$ and $b=q(A)$, so $a-b=d$.
The grouped bound is the Bernoulli KL $h_b(a)$.
For $a,b\in(0,1)$, $h_b(b)=h_b'(b)=0$ and
\[
 h_b''(a)=\frac1{a(1-a)}\ge4.
\]
Integrating twice between $b$ and $a$ gives $h_b(a)\ge2(a-b)^2=2d^2$.
Endpoint cases follow by limits or have infinite KL.
Thus $\KL(P\Vert q)\ge2\TV(P,q)^2$.
\end{proof}

\subsection{Limits of recovering distributions from means alone}
\begin{proposition}[Two-distribution recovery lower bounds]\label{prop:lower}
Suppose a recovery rule receives only the exact mean and outputs a distribution $Q(m)$.
Let $P_0,P_1$ be distinct distributions with the same mean, and set $d=\TV(P_0,P_1)$.
For $q=Q(m)$,
\begin{align}
 \max_{j\in\{0,1\}}\TV(P_j,q)&\ge d/2,\\
 \frac{\mathcal E_{P_0}(q)+\mathcal E_{P_1}(q)}2
 &\ge\frac{D(P_0,P_1)}{4(K-1)},\label{eq:twopointrps}\\
 \max_{j\in\{0,1\}}\KL(P_j\Vert q)&\ge d^2/2.
\end{align}
The first two bounds are sharp when the feasible forecast set includes $(P_0+P_1)/2$.
\end{proposition}
\begin{proof}
The rule sees the same input under both truths and therefore must use the same forecast.
The triangle inequality gives $d\le\TV(P_0,q)+\TV(q,P_1)$, proving the first bound.
Let $v_j=(F_{P_j}(k))_{k=0}^{K-2}$ and $u=(F_q(k))_{k=0}^{K-2}$.
The Euclidean identity
\[
 \frac{\|v_0-u\|^2+\|v_1-u\|^2}{2}
 =\left\|u-\frac{v_0+v_1}{2}\right\|^2+\frac{\|v_0-v_1\|^2}{4}
\]
and Proposition~\ref{prop:rps} give Eq.~\eqref{eq:twopointrps}.
The mixture forecast has the midpoint CDF and TVD $d/2$ to each truth, proving sharpness of those bounds.
For at least one truth the TVD is at least $d/2$; Proposition~\ref{prop:pinsker} then gives KL at least $2(d/2)^2$.
\end{proof}

For a randomized rule, the RPS bound holds after averaging over its randomization, and the corresponding TVD and KL worst-case bounds hold for expected losses.
Indeed the same inequalities can first be averaged over the common output law, using $\max_j\E L_j\ge(\E L_0+\E L_1)/2$ and convexity where necessary.
For raw expected RPS, the right side of Eq.~\eqref{eq:twopointrps} must be added to the average of the two irreducible scores in Eq.~\eqref{eq:riskrps}.
These are information bounds for a rule given only means.
They do not lower-bound an estimator that also receives informative individual ratings, item features, or valid shape restrictions.
Nor do they force every correlation optimizer to choose an inaccurate member of its optimum set.

\section{Consequences of model restrictions}
\subsection{Global and local identification}
Assume the model is well specified on the chosen items, with $P_i=q_{\theta^\ast}(\cdot\mid x_i)$ and nonconstant $t=M(\theta^\ast)$.
Proposition~\ref{prop:invariance} gives the exact global optimum set
\begin{equation}\label{eq:thetaset}
 \Theta_\rho=\{\theta\in\Theta:M(\theta)=a\mathbf1+bM(\theta^\ast),\ a\in\R,\ b>0\}.
\end{equation}
Correlation identifies the true predictions exactly when every $\theta\in\Theta_\rho$ satisfies $q_\theta(\cdot\mid x_i)=P_i$ for all $i$.
It identifies the parameter itself only when this set reduces to $\{\theta^\ast\}$, after any stipulated parameter equivalences are removed.
This requires excluding alternative feasible affine copies and alternative shapes with the same means.
Within a well-specified class, CE also identifies the true predictions, by Proposition~\ref{prop:ce}; parameter symmetries can still prevent parameter identification.

Local identification replaces $\Theta$ in this criterion by a neighborhood of $\theta^\ast$.
It does not exclude a distant second optimum.
For an open $d$-dimensional parameter domain and continuously differentiable $M$, the following is a sufficient local condition:
\begin{equation}\label{eq:tangent}
 DM(\theta^\ast)v\in\operatorname{span}\{\mathbf1,t\}\quad\Longrightarrow\quad v=0.
\end{equation}
To see this, define the standardized mean map
\[
 B(\theta)=\frac{M(\theta)-\overline{M(\theta)}\mathbf1}{\|M(\theta)-\overline{M(\theta)}\mathbf1\|}.
\]
For any perturbation $h$ of the mean vector, its derivative under centering and normalization vanishes exactly when $h\in\operatorname{span}\{\mathbf1,t\}$ at $\theta^\ast$.
Thus Eq.~\eqref{eq:tangent} makes $DB(\theta^\ast)$ have full column rank.
Selecting $d$ independent output coordinates and applying the inverse function theorem shows that $B$ is injective in a neighborhood.
Since correlation one is equivalent to $B(\theta)=B(\theta^\ast)$, the local optimum is unique there.
Failure of this first-order condition alone does not prove nonidentification; a singular derivative can coexist with local injectivity.

A family name, including logistic-linear, does not by itself establish Eq.~\eqref{eq:thetaset}'s uniqueness on a finite design.
The number and placement of the items and any shared-parameter symmetries matter.
A claim about identities of functions on an interval is stronger than a claim about equality on finitely many items.

Restricted models can also contain approximate affine copies.
For softmax logits $c v_i(r)$ with fixed $v_i$ and $c\to0$, direct Taylor expansion gives
\begin{equation}\label{eq:weaklogit}
 m_i(c)=\frac{K-1}{2}+\frac cK\sum_r\left(r-\frac{K-1}{2}\right)v_i(r)+O(c^2).
\end{equation}
Indeed $q_i(r;c)=K^{-1}+cK^{-1}(v_i(r)-\bar v_i)+O(c^2)$, and multiplication by $r$ and summation gives Eq.~\eqref{eq:weaklogit}.
On a fixed finite design the remainder is uniform across items.
If the displayed linear coefficient is nonconstant, then for any fixed $a>0$, $\rho(m(c),m(ac))\to1$ as $c\to0^+$.
This establishes approximate scale ambiguity near uniform logits, not an exact additional symmetry of every softmax model.
It supplies no empirical explanation for a fitted model without checking that model's regime.

\subsection{Linear mean models}
\begin{proposition}[Best linear direction]\label{prop:linear}
Let $\phi$ be a random feature vector with positive definite covariance $\Sigma$, let $t$ be a target mean with $\Var(t)=\sigma_t^2>0$, and put $c=\Cov(\phi,t)\ne0$.
For unconstrained linear mean predictions $m=a+w^\top\phi$,
\[
 \rho(w)=\frac{w^\top c}{\sqrt{w^\top\Sigma w}\,\sigma_t}.
\]
Its maximum is attained exactly at $w=\beta\Sigma^{-1}c$, $\beta>0$, with
\[
 \rho_{\max}^2=\frac{c^\top\Sigma^{-1}c}{\sigma_t^2}.
\]
The offset $a$ and positive scale $\beta$ are undetermined.
\end{proposition}
\begin{proof}
Centering removes $a$, and the covariance and variance of $w^\top\phi$ give the first formula.
Cauchy--Schwarz applied to $\Sigma^{1/2}w$ and $\Sigma^{-1/2}c$ gives
\[
 (w^\top c)^2\le(w^\top\Sigma w)(c^\top\Sigma^{-1}c).
\]
Equality requires proportionality of these vectors, and a positive correlation maximum requires a positive proportionality factor.
This yields $w=\beta\Sigma^{-1}c$.
The same vector solves the normal equations for least squares with an intercept.
Its explained variance is $c^\top\Sigma^{-1}c$, giving the stated squared correlation.
\end{proof}

The result also applies to empirical covariances on a fixed item set.
If $c=0$, every nonconstant feasible linear forecast has correlation zero, so there is no distinguished direction.
If forecasts must lie in $[0,K-1]$, the optimum above requires a feasible offset and scale; the unconstrained formula alone does not guarantee feasibility on unbounded feature supports.
Correlation can equal one in this linear class only when the target is an affine function of the features almost surely.
A linear mean head itself provides no distributional shape output.

For comparison, a softmax-linear distribution model $q_\theta(r\mid x)=\softmax(W\phi(x))_r$ has expected CE gradient
\[
 \frac{\partial L}{\partial W_r}=\E\big[\phi(x)(q_\theta(r\mid x)-P(r\mid x))\big].
\]
This follows by applying Eq.~\eqref{eq:gradients} and the chain rule.
A finite unpenalized stationary point matches these feature-weighted category frequencies.
Under misspecification this need not recover the true conditional distribution or each item's mean.
The moment-matching guarantee in Proposition~\ref{prop:exponential} assumes a free scalar natural parameter for the distribution being fitted; arbitrary shared-parameter classes do not inherit that guarantee.

\subsection{Nonlinear mean models}
A nonlinear mean model can represent a nonlinear target $t(x)$ and attain population correlation one.
Its input-to-rating map need not be linear.
If an additional output changes the distribution along a feasible path while preserving all means, the objective is constant on that path.
If shared parameters force every shape change to alter means, correlation can influence shape indirectly, as Proposition~\ref{prop:gradient} states.
Thus the unidentifiable set must be intersected with the feasible model class rather than inferred from a parameter's informal description as a spread parameter.

\section{Finite ratings and estimation error}
\subsection{Noise in the observed target mean}
Let $I$ be a random held-out item and $Y_{I1},\ldots,Y_{Ik}$ be conditionally independent ratings drawn from $P_I$.
Set $t_I=\E[Y\mid I]$, $\sigma_I^2=\Var(Y\mid I)$, and $\widehat t_I=k^{-1}\sum_{j=1}^kY_{Ij}=t_I+\varepsilon_I$.
Then
\[
 \E[\varepsilon_I\mid I]=0,\qquad \Var(\varepsilon_I\mid I)=\frac{\sigma_I^2}{k}.
\]
For a fixed forecast mean $m_I$, or a forecast conditioned on training data independent of these test ratings, the law of total covariance yields
\[
 \Cov(m_I,\varepsilon_I)=\E[(m_I-\E m_I)\E(\varepsilon_I\mid I)]=0.
\]
The same argument gives $\Cov(t_I,\varepsilon_I)=0$.
The law of total variance therefore implies
\[
 \Cov(m_I,\widehat t_I)=\Cov(m_I,t_I),\qquad
 \Var(\widehat t_I)=\Var(t_I)+\E[\sigma_I^2]/k.
\]
Whenever the relevant variances are positive,
\begin{equation}\label{eq:attenuation}
 \operatorname{Corr}(m_I,\widehat t_I)
 =\frac{\operatorname{Corr}(m_I,t_I)}{\sqrt{1+\E[\sigma_I^2]/(k\Var(t_I))}}.
\end{equation}
Conditional mean-zero noise suffices; independence of the noise from the item is unnecessary and generally false with heteroskedastic ratings.
For item-dependent counts $k_I$, replace $\E[\sigma_I^2]/k$ by $\E[\sigma_I^2/k_I]$ under the same sampling assumptions.

Equation~\eqref{eq:attenuation} is a population correlation identity over random held-out items and fresh ratings.
It is not an equality for the expectation of a finite-sample Pearson coefficient, whose denominator is random.
It also does not apply directly to a predictor fitted on the same rating noise.
An unrestricted predictor can interpolate $m_i=a+b\widehat t_i$ whenever feasible and the observed means are nonconstant, attaining training correlation one while memorizing noise.
Correlated raters, changing rater populations, and selection of which items receive ratings require a different noise analysis.

\subsection{A regular asymptotic CE calculation}
The familiar $P/(2N)$ expression concerns a specific estimation regime, not every distributional learner.
To state it precisely, let $Z_j=(X_j,Y_j)$, $j=1,\ldots,N$, be independent and identically distributed observations from a well-specified parametric conditional model $q_{\theta^\ast}(Y\mid X)$.
Let $\theta^\ast$ lie in the interior of an identifiable $P$-dimensional parameter domain.
Assume a twice differentiable population CE risk
\[
 R(\theta)=\E[-\log q_\theta(Y\mid X)]
\]
with nonsingular Fisher information $I=\nabla^2R(\theta^\ast)$, and an efficient consistent unpenalized maximum-likelihood estimator satisfying
\[
 \sqrt N(\widehat\theta-\theta^\ast)\ \xrightarrow{d}\ \mathcal N(0,I^{-1}).
\]
Assume in addition the moment convergence needed to take expectations of the quadratic form and an expected Taylor remainder $o(N^{-1})$.
For example, appropriate uniform integrability of $N\|\widehat\theta-\theta^\ast\|^2$ together with control of the remainder suffices.
Then
\begin{align}
 R(\widehat\theta)-R(\theta^\ast)
 &=\tfrac12(\widehat\theta-\theta^\ast)^\top I(\widehat\theta-\theta^\ast)+\text{remainder},\\
 \E[R(\widehat\theta)-R(\theta^\ast)]
 &=\frac1{2N}\operatorname{tr}(II^{-1})+o(N^{-1})
 =\frac{P}{2N}+o(N^{-1}).\label{eq:asymptotic}
\end{align}
The first-order Taylor term vanishes at the interior population minimizer.
The second equality uses $\E[N(\widehat\theta-\theta^\ast)(\widehat\theta-\theta^\ast)^\top]\to I^{-1}$ and the trace identity for quadratic forms.
Thus the expectation statement requires more than convergence in distribution alone.
The result is a conditional regular-asymptotic calculation, measured in nats per fresh observation.

When $N$ is comparable to $P$, when the MLE lies near a boundary or does not exist, under misspecification, or with strong regularization, these assumptions can fail.
Overparameterized models and clustered ratings need separate analyses.
In particular, an estimated zero probability can yield infinite test CE even when the corresponding category was absent from training.
A probability floor changes the estimator and its evaluated loss; it is not implicit in the preceding identities.

\subsection{Why there is no universal break-even sample size}
Correlation's population optimum set includes the truth when the class is well specified.
Its nonuniqueness therefore does not force a positive error floor for every fitting procedure.
Initialization, regularization, a tie rule, and model restrictions can select different members of that set.
If a specified procedure converges to a forecast with positive excess CE $B$, then comparing $P/(2N)$ to $B$ suggests the heuristic scale $N\approx P/(2B)$.
This comparison requires both risks to use CE in nats and the asymptotic assumptions to be relevant.
It is invalid to substitute an RPS gap for $B$ in that expression.
Neither $B>0$ nor a universal crossover follows from correlation's invariances.

Properness describes the optimizer of population expected loss.
It does not order finite-sample estimators, optimization methods, or validation procedures.
With a small number of ratings, a restricted or regularized forecast can have lower held-out RPS or CE than a more flexible fit.
The invariance results also do not compare choosing among existing judges with fitting a new distributional head.
These are distinct statistical procedures whose performance depends on the candidate class and data.

\section{An exact reproducible illustration}\label{sec:example}
Consider three items on $R=\{0,1,2\}$ with true and forecast distributions
\begin{equation}\label{eq:example}
 P=\begin{pmatrix}
 3/5&3/10&1/10\\
 1/5&3/5&1/5\\
 1/10&3/10&3/5
 \end{pmatrix},\qquad
 Q=\begin{pmatrix}
 13/20&1/5&3/20\\
 2/5&1/5&2/5\\
 3/20&1/5&13/20
 \end{pmatrix}.
\end{equation}
Rows index items and columns index rating categories in increasing order.
All probabilities are strictly positive.
The true and forecast means both equal $t=m=(1/2,1,3/2)$, so their Pearson correlation across items is exactly one.
Nevertheless, $Q_i-P_i=s_i(1,-2,1)$ with $(s_1,s_2,s_3)=(1/20,1/5,1/20)$.
Thus $\TV(P_i,Q_i)=2s_i$ and the two CDF differences equal $(-s_i,s_i)$.
Proposition~\ref{prop:rps} gives $\mathcal E_{P_i}(Q_i)=s_i^2$.
Consequently,
\[
 \frac13\sum_i\TV(P_i,Q_i)=\frac15,\qquad
 \frac13\sum_i\mathcal E_{P_i}(Q_i)=\frac{(1/20)^2+(1/5)^2+(1/20)^2}{3}=\frac3{200}.
\]
The expected normalized RPS of the true distributions is $49/300$.
Adding the excess gives $107/600$ for $Q$.
These expected scores include stochastic rating noise; their difference isolates distributional error.

% Provenance: experiments/exact_counterexample/reproduce.py; exact rational inputs, standard-library logarithms for CE.
\begin{table}[ht]
\centering
\begin{tabular}{@{}lrrrr@{}}
\toprule
Forecast & Pearson & TVD & Excess RPS & CE (nats)\\
\midrule
True distribution $P$ & $1$ & $0$ & $0$ & $0.915387330$\\
Changed shape $Q$ & $1$ & $1/5$ & $3/200$ & $1.064735095$\\
\bottomrule
\end{tabular}
\caption{Exact three-item illustration, not a fitted-model comparison.
Pearson correlation uses the three conditional means; higher is better for association.
TVD compares each forecast with the true distribution and lies in $[0,1]$.
Excess normalized individual-rating RPS subtracts the true forecast's expected score and lies in $[0,1]$.
CE is expected individual-rating cross-entropy in nats.
All distributional metrics are averaged equally across the three items, and lower is better.
The true distribution is the reference forecast.
There is no train, validation, or test split, fitted judge, optimizer, or sampling uncertainty; every entry follows from the displayed distributions.}
\label{tab:exact}
\end{table}

The script \texttt{experiments/exact\_counterexample/reproduce.py} reproduces the table with Python's standard library.
It uses rational arithmetic for means, CDFs, TVD, and RPS and evaluates logarithms numerically only for CE.
It also verifies mass, positivity, the signed CDF identity, and equality of the mean vectors.
The average excess CE is approximately $0.149347765$ nats, which is the difference between the two CE values before rounding.
The example demonstrates nonidentification, not the behavior of an optimizer or the frequency of this error in practice.

\section{Limitations and implications}
The results describe population identification on a finite ordered rating scale.
They assume a specified conditional target distribution for each item.
In practice, a small number of ratings estimates that distribution imprecisely.
Changes in the rater population or the meaning of rating categories change the target itself.
One observed rating supports an unbiased individual-rating CE or RPS observation under the assumed sampling law, but it does not reveal the full conditional distribution for that item.

The geometry of unrestricted distributions does not automatically describe a restricted neural model.
Parameter sharing, regularization, initialization, and optimization may choose a particular shape even when the objective does not identify it.
Conversely, strict properness does not remove misspecification or optimization error.
The linear-model result needs covariance and feasibility conditions, and the exponential-family result needs its stated sufficient statistic and positive variance.
Global identification requires checking the full feasible affine-copy set; local derivative conditions provide only the stated local guarantee.

The exact illustration is a mathematical construction with no trained judge or sampled human ratings.
It establishes possibility and reproduces a loss difference exactly.
It does not estimate the prevalence of such differences or compare real-world training methods.
The asymptotic CE calculation concerns regular, interior, well-specified iid maximum likelihood with additional expectation control.
It does not predict a universal sample size at which CE-trained judges surpass correlation-trained judges.

Our main results concern Pearson correlation of conditional means.
Spearman and Kendall coefficients computed only from those means also discard equal-mean shape differences, but their equality conditions and gradients differ.
Probability-profile correlation is a separate objective with the ambiguity in Proposition~\ref{prop:profiles}.
For applications that require only rank association, distribution recovery can exceed the task's requirements.
For applications that interpret a judge's output as a rating distribution, individual-rating proper scores directly target that distribution.
Correlation remains a useful association diagnostic when its target and undefined constant-vector cases are reported explicitly.

\begin{thebibliography}{9}
\bibitem{gneiting2007}
Tilmann Gneiting and Adrian E. Raftery.
Strictly proper scoring rules, prediction, and estimation.
\emph{Journal of the American Statistical Association}, 102(477):359--378, 2007.
\href{https://doi.org/10.1198/016214506000001437}{doi:10.1198/016214506000001437}.

\bibitem{epstein1969}
Edward S. Epstein.
A scoring system for probability forecasts of ranked categories.
\emph{Journal of Applied Meteorology}, 8(6):985--987, 1969.
\href{https://doi.org/10.1175/1520-0450(1969)008\%3C0985:ASSFPF\%3E2.0.CO;2}{doi:10.1175/1520-0450(1969)008\textless0985:ASSFPF\textgreater2.0.CO;2}.
\end{thebibliography}
\bigskip
\noindent\small\copyright\ 2026 Brando Miranda.
This preprint is licensed under the \href{https://creativecommons.org/licenses/by/4.0/}{Creative Commons Attribution 4.0 International License (CC BY 4.0)}.

\end{document}
